\begin{exercise}{140}
Prove or disprove that the diagram
\[
\begin{tikzcd}[column sep=large]
 \mathrm{Fun}(\Delta^1,\mathcal C) \arrow[r] \arrow[d]
   & \mathrm{Fun}(\Delta^1,\mathcal D) \arrow[d] \\
 \mathcal C\times\mathcal C \arrow[r]
   & \mathcal D\times\mathcal D
\end{tikzcd}
\]
is homotopy-cartesian if and only if $f$ is fully faithful.

\begin{proof}
The assertion is true. Let $f\colon\mathcal C\to\mathcal D$ be the
functor inducing the horizontal maps.
Write $A(X)=\mathrm{Fun}(\Delta^1,X)$ and let
$p_X\colon A(X)\to X\times X$ be endpoint evaluation.
Proposition~2.2.5 makes every $p_X$ an isofibration. Put
\[
 \begin{gathered}
 P=(\mathcal C\times\mathcal C)
       \times_{\mathcal D\times\mathcal D}A(\mathcal D),\\
 q=(p_{\mathcal C},A(f))\colon A(\mathcal C)\longrightarrow P.
 \end{gathered}
\]
The projection $P\to\mathcal C\times\mathcal C$ is also an
isofibration. By Definition~3.1.9, the square is homotopy-cartesian
exactly when $q$ is a Joyal equivalence.
If so, Lemma~2.5.32 makes its fibre map over every $(x,y)$,
\[
 \mathrm{map}_{\mathcal C}(x,y)\longrightarrow
 \mathrm{map}_{\mathcal D}(fx,fy),
\]
a Joyal equivalence between Kan complexes, hence a homotopy
equivalence. Thus $f$ is fully faithful.

Conversely, suppose $f$ is fully faithful, and let
$\mathcal I\subseteq\mathcal D$ be the full subcategory spanned by
$f(\mathcal C_0)$.
The arrow mapping-space model gives
$\mathrm{map}_{\mathcal I}(u,v)=\mathrm{map}_{\mathcal D}(u,v)$
for $u,v\in\mathcal I_0$.
Thus the factor $e\colon\mathcal C\to\mathcal I$ is fully faithful
and surjective on objects, so Theorem~2.3.20 makes it a Joyal
equivalence.
There is a strict pullback identity
\[
 A(\mathcal I)=(\mathcal I\times\mathcal I)
       \times_{\mathcal D\times\mathcal D}A(\mathcal D):
\]
indeed, all vertices of $\Delta^n\times\Delta^1$ lie on its endpoint
restrictions, so an arrow simplex with endpoints in $\mathcal I$
lies entirely in this full subcategory.
Consequently,
\[
 P\cong(\mathcal C\times\mathcal C)
       \times_{\mathcal I\times\mathcal I}A(\mathcal I).
\]
Exercises~\ref{ex:90} and~\ref{ex:101} make $A(e)$ and $e\times e$
Joyal equivalences.
The projection $r\colon P\to A(\mathcal I)$ is the pullback of
$e\times e$ along the isofibration $p_{\mathcal I}$, so it is a
Joyal equivalence by Proposition~2.5.5.
Since $rq=A(e)$, two-out-of-three makes $q$ a Joyal equivalence,
as required.
\end{proof}
\end{exercise}
